\begin{lemma}[Lemma 3.2]
Let $n\ge2$, $p>1$ and $r>0$, and suppose the Sobolev--Poincar\'e inequality on balls holds in dimension $n$
(\noderef{CitedSobolevPoincareBall}).  Then there exist constants
$\theta_2=\theta_2(n,p,r)\in(0,1)$ and $\nu=\nu(n,p,r)\in(0,1)$ with the
following property.  Let $\varepsilon>0$, $\gamma_0>0$, $\gamma_1\le r\gamma_0$,
let $x_0\in\R^n$ and $R>0$, write $B_R=B_R(x_0)$,
let $A$ satisfy (Q1) of the thesis--(Q2) of the thesis on $B_R$ with parameters
$p,\varepsilon,\gamma_0,\gamma_1$ (\noderef{EllipticStructureQ12}), assume the
cited regularity result (C1) for $\varepsilon>0$
(\noderef{CitedRegularityEpsPos}), let $u$ be a
weak solution of $\operatorname{div}A(\nabla u)=0$ in $B_R$
(\noderef{IsWeakSolutionOn}) with weak gradient $g=\nabla u$, suppose that for
each $i\in\{1,\dots,n\}$ the component $g^i=D_iu$ belongs to
$W^{1,2}_{\mathrm{loc}}(B_R)$ with weak gradient $Dg^i$
(\noderef{MemW1pLocOn}) --- that is, $u$ has the $W^{2,2}_{\mathrm{loc}}$
regularity that the cited result (C1) supplies for $\varepsilon>0$, the exponent
being $2$ regardless of $p$ --- and let
\[
\mu=\max_{1\le k\le n}\|D_ku\|_{L^\infty(B_R)} ,
\]
i.e.\ $\mu$ is the least number with $|D_ku|\le\mu$ a.e.\ in $B_R$ for every
$k$.  Then:
\begin{itemize}
\item[\textup{(a)}] For every $k$, $1\le k\le n$: if
\[
\bigl|\{x\in B_R:D_ku(x)>\tfrac12\mu\}\bigr|\ge\theta_2|B_R|
\quad\text{or}\quad
\bigl|\{x\in B_R:D_ku(x)<-\tfrac12\mu\}\bigr|\ge\theta_2|B_R| ,
\]
then
\[
|D_ku(x)|\ \ge\ \min\{1/16,\ 2^{-1-4/p}\}\cdot\mu\qquad\text{for a.e.\ }x\in B_{R/2}.
\]
The index $k$ is shared between the hypothesis and the conclusion, as in the
thesis: the nondegeneracy of one particular partial derivative on a set of
measure at least $\theta_2|B_R|$ forces that same partial derivative to stay
away from $0$ on the half ball.
\item[\textup{(b)}] If the inequalities of \textup{(a)} hold for no $k$,
$1\le k\le n$, then
\[
\max_{1\le k\le n}\|D_ku\|_{L^\infty(B_{R/2})}\le\nu\cdot\mu ,
\]
i.e.\ $|D_ku|\le\nu\mu$ a.e.\ in $B_{R/2}$ for every $k$.
\end{itemize}

The hypothesis \noderef{CitedSobolevPoincareBall} is the cited result (C6) of the
paper's Note for the formalization, carried here as an explicit argument because
the proof uses it --- it enters through the constants of \noderef{LevelTruncationMaximumPrinciple} and \noderef{LevelSetMeasureImprovement} --- which is what the
Note permits for the results the thesis quotes without proof.  Although it is
bound before the existential for $\theta_2$ and $\nu$, it introduces no dependence of
$\theta_2$ and $\nu$ on any quantity outside the list the paper names: (C6) is a statement
about $n$ alone, and $n$ is already bound first, so the hypothesis carries no
numerical content of its own.
\end{lemma}
\begin{proof}
Throughout, $B_\rho=B_\rho(x_0)$, $w=\varepsilon+|\nabla u|^2$, $g=\nabla u$,
$g^k=D_ku$ and $Dg^k$ is the weak gradient of $g^k$ supplied by the hypothesis
(\noderef{MemW1pLocOn}).  All constants called $C$ with a subscript are positive
and depend only on $n$, $p$ and $r$; this is checked at each occurrence.  We
write $\omega_n=|B_1|$, so $|B_\rho|=\omega_n\rho^{\,n}$ (Preamble: the volume of
a ball).

\emph{Step 0: normalisations, and the five facts used repeatedly.}

\textbf{(0a) The degenerate case $\mu=0$.}  If $\mu=0$ then $|g^k|\le0$ a.e.\ on
$B_R$ for every $k$, so $g^k=0$ a.e.\ there.  Part (a)'s conclusion reads
$0\le|g^k|$ a.e.\ on $B_{R/2}$, which holds, and part (b)'s reads $|g^k|\le0$
a.e.\ on $B_{R/2}$, which holds as well.  Both parts are therefore true for any
choice of $\theta_2,\nu\in(0,1)$, and we may and do assume $\mu>0$ from now on.
Put
\[
\lambda=\frac\mu2>0,\qquad h_0=\lambda^{p/2}>0 .
\]
Note $\mu\ge0$ always, because $|g^k|\le\mu$ a.e.\ on the set $B_R$ of positive
measure.

\textbf{(0b) a continuous representative $g'$ of the gradient, and $w$ is
pinched.}  The cited regularity result
(C1) (\noderef{CitedRegularityEpsPos}), applied with the present
$\varepsilon>0$, the structure (Q1) of the thesis--(Q2) of the thesis
(\noderef{EllipticStructureQ12}) and the weak solution $u$
(\noderef{IsWeakSolutionOn}), gives among its conclusions a field $g'$ with
$g'$ continuous on $B_R$ and $g=g'$ almost everywhere on $B_R$.  \emph{It does
not assert continuity of the given $g$, and cannot: a weak gradient is an
almost-everywhere class, so continuity is a property of a representative.  The
argument below is therefore run with $g'$, and the pinch hypothesis is
transported to it across the almost-everywhere identity.}

First, $|g'^k|\le\mu$ \emph{everywhere} on $B_R$.  The hypothesis gives
$|g^k|\le\mu$ almost everywhere on $B_R$, and $g=g'$ almost everywhere there, so
$|g'^k|\le\mu$ almost everywhere on $B_R$; the exceptional set
$\{x\in B_R:|g'^k(x)|>\mu\}$ is open, because $g'$ is continuous on the open set
$B_R$, and it is null by the previous sentence, and a
nonempty open subset of $\R^n$ contains a ball of positive measure (Preamble: the
volume of a ball), so it is empty.  Put
\[
w'(x)=\varepsilon+|g'(x)|^2 ,
\]
so that $w'=w$ almost everywhere on $B_R$ and, by the previous paragraph,
everywhere on $B_R$,
\[
\varepsilon\ \le\ w'\ =\ \varepsilon+\sum_{i}(g'^i)^2\ \le\ \varepsilon+4n\lambda^2 ,
\tag{0.1}
\]
using $|g'^i|\le\mu=2\lambda$ for each of the $n$ components.  The two-sided
bound (0.1) therefore holds for $w$ almost everywhere on $B_R$ as well, which is
the form in which the steps below that read $w$ under an integral consume it.

\textbf{(0c) $\gamma_0\le\gamma_1$.}  The clauses \texttt{Q1} and \texttt{Q2} of
\noderef{EllipticStructureQ12} each hold for almost every $x$, so they hold
simultaneously at almost every $x$, and $B_R$ has positive measure; fix such an
$x\in B_R$.  Take $z$ arbitrary, $h=0$ and $\xi=e_1$: the guard
$\varepsilon+\|0\|^2=\varepsilon\ne0$ holds because $\varepsilon>0$.  Then
\texttt{Q1} gives
$\gamma_0\varepsilon^{p/2-1}\le\sum_i(D A(0)e_1)_i(e_1)_i=(DA(0)e_1)_1$ and
\texttt{Q2} gives $|(DA(0)e_1)_1|\le\gamma_1\varepsilon^{p/2-1}$, whence
$\gamma_0\varepsilon^{p/2-1}\le\gamma_1\varepsilon^{p/2-1}$ and, dividing by
$\varepsilon^{p/2-1}>0$, $\gamma_0\le\gamma_1$.  (This is needed because
\noderef{DifferentiatedEquationWeak} carries it as a hypothesis, while the
present lemma carries $\gamma_1\le r\gamma_0$ instead.)

\textbf{(0d) Removing a pinched weight.}  The passage from a weighted
Caccioppoli inequality to an instance of (2.1) of the thesis is
\noderef{PinchedCaccioppoliLevelBound}, applied in Step 3; what this proof owes
it is the level-set pinch of Step 2 --- the paper's own form at l.~884--885 ---
and no case distinction on the sign of $p/2-1$ is needed, that being a property
of that node's constant and not of this proof.

\textbf{(0e) The cut-off.}  Let $\rho_{1/8}$ be the standard mollifier of the
Preamble in dimension one at scale $\tfrac18$ --- smooth, nonnegative, supported
in $[-\tfrac18,\tfrac18]$, of total integral $1$ --- and put
$\psi_\flat=\mathbf1_{(-\infty,1/2]}*\rho_{1/8}$.  Then
$\psi_\flat\in C^\infty(\R)$ with
$\psi_\flat'=\mathbf1_{(-\infty,1/2]}*\rho_{1/8}'$ (Preamble: convolution with a
smooth compactly supported kernel, the indicator being locally integrable),
$0\le\psi_\flat\le1$, $\psi_\flat=1$ on $(-\infty,\tfrac38]$ and $\psi_\flat=0$
on $[\tfrac58,\infty)$ --- for $t\le\tfrac38$ every $s$ with $|s|\le\tfrac18$ has
$t-s\le\tfrac12$, so the convolution integrates $\rho_{1/8}$ over its whole
support, and for $t\ge\tfrac58$ every such $s$ has $t-s\ge\tfrac12$, so the
integrand vanishes almost everywhere --- and
$|\psi_\flat'|\le\|\rho_{1/8}'\|_{L^1(\R)}=:\Lambda<\infty$, an absolute
constant.  Given $0<\rho_0<\rho_1$ put
\[
\zeta_{\rho_0,\rho_1}(x)=\psi_\flat\Bigl(\frac{|x-x_0|-\rho_0}{\rho_1-\rho_0}\Bigr).
\]
Then $0\le\zeta\le1$; $\zeta=1$ on $\overline B_{\rho_0}$; $\zeta=0$ outside
$\overline B_{\rho_0+\frac58(\rho_1-\rho_0)}$, so $\operatorname{supp}\zeta$ is
compact and contained in $B_{\rho_1}$; $\zeta$ is smooth, because
$x\mapsto|x-x_0|$ is smooth away from $x_0$ and $\zeta$ is constant on the
neighbourhood $B_{\rho_0}$ of $x_0$; and
$|\nabla\zeta|\le\Lambda/(\rho_1-\rho_0)$, with $\nabla\zeta$ supported in
$B_{\rho_1}$.  (This cut-off is built inline from the Preamble's mollifier
rather than cited from a node, as the run's refuted \texttt{StandardCutoff}
proposal requires.)

\textbf{(0f) Localisation.}  For $0<T<R$ one has
$\operatorname{closure}B_T=\overline B_T$ compact and $\overline B_T\subseteq
B_R$, so clause (4) of \noderef{WeakGradientRestrict} (with $q=2$, the open sets
$B_R$ and $B_T$) upgrades the hypothesis
$g^k\in W^{1,2}_{\mathrm{loc}}(B_R)$ to
\[
g^k\in W^{1,2}(B_T)\ \text{with weak gradient }Dg^k
\qquad(0<T<R).
\tag{0.3}
\]
Conversely, if a pair $(f,G)$ satisfies $f\in W^{1,2}(B_T)$ with weak gradient
$G$ for every $T<R$, then $f\in W^{1,2}_{\mathrm{loc}}(B_R)$ with weak gradient
$G$: the defining identity of \noderef{HasWeakGradientOn} on $B_R$ quantifies
over test functions $\varphi$ with $\operatorname{tsupport}\varphi\subseteq B_R$
compact, and every such $\operatorname{tsupport}\varphi$ lies in some $B_T$ with
$T<R$, where the identity is available; and every compact $K\subseteq B_R$ lies
in such a $B_T$, so the $L^2$ bounds of \noderef{MemW1pLocOn} follow from those
of \noderef{MemW1pOn} on $B_T$.  We refer to this passage as (0.3$'$).  Finally
$\operatorname{volume}B_T\ne\infty$ for every $T$ (Preamble: the volume of a
ball), which is the finiteness hypothesis of
\noderef{WeakGradientPositivePart}, \noderef{WeakGradientC1Comp} and clause (1)
of \noderef{WeakGradientSmoothMul}.

\emph{Step 1: inequality \textup{(3.4)}.}  Apply
\noderef{CaccioppoliLevelWeight} with $n,p,r$ and let $c_4>0$ be the constant it
produces; it depends only on $n,p,r$.  Its remaining hypotheses are the present
lemma's $0<\varepsilon$, $0<\gamma_0$, $\gamma_1\le r\gamma_0$, the openness of
$B_R$, (Q1) of the thesis--(Q2) of the thesis (\noderef{EllipticStructureQ12}), the cited
regularity (C1) (\noderef{CitedRegularityEpsPos}), the weak solution property
(\noderef{IsWeakSolutionOn}) and
$g^i\in W^{1,2}_{\mathrm{loc}}(B_R)$ for every $i$
(\noderef{MemW1pLocOn}).  Hence for every $k$, every $\ell\in\R$ and every
$\zeta$ smooth with compact support in $B_R$,
\[
\int_{B_R\cap\{g^k>\ell\}}w^{\frac p2-1}\|Dg^k\|^2\zeta^2\,dx
\ \le\ c_4\int_{B_R}w^{\frac p2-1}\bigl((g^k-\ell)^+\bigr)^2|\nabla\zeta|^2\,dx
\tag{3.4$^+$}
\]
and
\[
\int_{B_R\cap\{g^k<\ell\}}w^{\frac p2-1}\|Dg^k\|^2\zeta^2\,dx
\ \le\ c_4\int_{B_R}w^{\frac p2-1}\bigl((\ell-g^k)^+\bigr)^2|\nabla\zeta|^2\,dx .
\tag{3.4$^-$}
\]
These are the two conclusions of that node verbatim, with
$\sum_i(\partial_i\zeta)^2=|\nabla\zeta|^2$ and
$\max(g^k-\ell,0)=(g^k-\ell)^+$.

\emph{Step 2: the pinch data for the two regimes \textup{(3.9)} and
\textup{(3.23)}.}  Both regimes feed the same node, and what each has to produce
is a pair $\kappa_1,\kappa_2$ with $0<\kappa_1\le w\le\kappa_2\le K\kappa_1$
almost everywhere on the set over which the corresponding conclusion integrates,
with $K=1+4n$ in both cases.  Apply \noderef{PinchedCaccioppoliLevelBound} with
$n$, the exponent $s=p/2-1$, the ratio $K=1+4n\ge1$ and the Caccioppoli constant
$c=c_4>0$ of Step 1, and let
\[
\gamma_9>0
\]
be the constant $\gamma$ it produces; it depends only on $n$, $s=p/2-1$,
$K=1+4n$ and $c_4$, hence only on $n,p,r$.  Two instantiations are used.

\textup{(i)} \emph{If $\varepsilon\ge\lambda^2$}, then by (0.1)
$\varepsilon\le w\le\varepsilon+4n\lambda^2\le(1+4n)\varepsilon$ on all of
$B_R$, so $\kappa_1=\varepsilon>0$, $\kappa_2=(1+4n)\varepsilon$ and $K=1+4n$ are
admissible for every level $\ell$ --- and, the bound holding on all of $B_R$, on
each of the two level sets $B_R\cap\{g^k>\ell\}$ and $B_R\cap\{g^k<\ell\}$
simultaneously, so that both of the node's hypotheses (H3$^+$) and (H3$^-$) are
available at every $\ell\in\R$.  This is the paper's (3.9) of the thesis.  (Only the
$(2.1)$ form is needed here; the paper also records that (3.9) of the thesis is of the
form $(2.11)$, which is used in Theorem~Theorem~2 of the thesis and not below.)

\textup{(ii)} \emph{If $\varepsilon<\lambda^2$ and $|\ell|\in[\lambda,2\lambda]$},
then on $\{g^k>\ell\}$ with $\ell\in[\lambda,2\lambda]$ we have
$g^k>\ell\ge\lambda$, so $w\ge(g^k)^2\ge\lambda^2$; and by (0.1)
$w\le\varepsilon+4n\lambda^2<(1+4n)\lambda^2$.  Thus $\kappa_1=\lambda^2>0$,
$\kappa_2=(1+4n)\lambda^2$, $K=1+4n$ are admissible \emph{on the level set},
i.e.\ the node's hypothesis (H3$^+$) holds for every $\ell\in[\lambda,2\lambda]$,
that is for every $\ell\in[\mu/2,\mu]$.  Symmetrically, on $\{g^k<\ell\}$ with
$\ell\in[-2\lambda,-\lambda]$ we have $g^k<\ell\le-\lambda$, so again
$w\ge\lambda^2$, and (H3$^-$) holds for every $\ell\in[-\mu,-\mu/2]$.  This is
the paper's (3.23) of the thesis, and it is here that the level-set form of the pinch
is indispensable: on all of $B_R$ the lower bound for $w$ is only $\varepsilon$,
which is smaller than $\lambda^2$ in this case by assumption, so the ratio
$\kappa_2/\kappa_1$ would not be bounded in terms of $n$ alone.  It is also here
that the \emph{decoupling} of the node's two implications is indispensable: at a
level $\ell\in[\mu/2,\mu]$ the $-$ pinch is simply false, since $g^k$ may be
arbitrarily close to $0$ on $\{g^k<\ell\}$, so a form with both pinches as common
hypotheses could not be applied in this regime at all.

\emph{Step 3: the instances of \textup{(2.1)}.}  Apply
\noderef{PinchedCaccioppoliLevelBound} --- the instance of Step 2, with the
constant $\gamma_9$ --- at the centre $x_0$, the radius $R$, the bound
$\Theta=\mu$, a level $\ell$, the pinch constants $\kappa_1,\kappa_2$ of
Step 2(i) resp.\ 2(ii), the function $f=g^k$, the field $G=Dg^k$ and the weight
$\widetilde w$ of the next paragraph.  Its hypotheses are: $0<R$; $0<\kappa_1$
and $\kappa_2\le K\kappa_1$ with $K=1+4n$, which is Step 2;
$\widetilde w$ measurable; $g^k\in W^{1,2}_{\mathrm{loc}}(B_R)$ with
weak gradient $Dg^k$ (\noderef{MemW1pLocOn}), which is the present lemma's
hypothesis; and $|g^k|\le\mu$ a.e.\ on $B_R$, which is (0b).

\emph{The weight is passed as a globally measurable modification.}
\noderef{PinchedCaccioppoliLevelBound} asks for a weight that is measurable on
all of $\R^n$, while (0b) gives a representative $g'$ continuous only on $B_R$.
Put
\[
\widetilde w(x)=\varepsilon+|g'(x)|^2=w'(x)\ \text{ for }x\in B_R,
\qquad
\widetilde w(x)=\varepsilon\ \text{ for }x\notin B_R .
\]
Then $\widetilde w$ is measurable on $\R^n$: $B_R$ is open, hence measurable,
and $\{B_R,\R^n\setminus B_R\}$ is a measurable partition on each piece of which
$\widetilde w$ is measurable --- on $B_R$ because $g'$ is continuous there by
(0b), so $x\mapsto\varepsilon+|g'(x)|^2$ is continuous on that open set,
and on the complement because $\widetilde w$ is constant there.  And
$\widetilde w=w'$ on $B_R$, hence $\widetilde w=w$ almost everywhere on $B_R$,
which is all the consumer below needs since every hypothesis and conclusion of
\noderef{PinchedCaccioppoliLevelBound} that reads the weight reads it under an
integral or almost everywhere.  Every hypothesis and every conclusion of
\noderef{PinchedCaccioppoliLevelBound} reads its weight only inside the ball ---
its two pinch hypotheses are almost-everywhere statements on
$B_R\cap\{f>\ell\}$ resp.\ $B_R\cap\{f<\ell\}$, its two Caccioppoli hypotheses
integrate over $B_R$ and subsets of it, and its conclusions do not mention the
weight at all --- so applying it with $\widetilde w$ gives exactly the
conclusions below, and Step 2's pinch bounds and Step 1's inequalities, which
are statements about $w$ on $B_R$, are statements about $\widetilde w$ verbatim.

Its two implications then read, for all $0<\rho_0<\rho_1<R$:
\begin{itemize}
\item given (H3$^+$) --- the pinch on $B_R\cap\{g^k>\ell\}$ --- and (H4$^+$) ---
  which is (3.4$^+$) of Step 1 at the level $\ell$, with $s=p/2-1$ and $c=c_4$,
  verbatim ---
  $\mathrm{DeGiorgiLevelBound}\ n\ \gamma_9\ g^k\ Dg^k\ x_0\ \ell\ \rho_0\ \rho_1$;
\item given (H3$^-$) and (H4$^-$) --- which is (3.4$^-$) ---
  $\mathrm{DeGiorgiLevelBound}\ n\ \gamma_9\ (-g^k)\ (-Dg^k)\ x_0\ (-\ell)\
  \rho_0\ \rho_1$.
\end{itemize}
So, by Step 2(i), in the regime $\varepsilon\ge\lambda^2$ both instances are
available at \emph{every} level $\ell\in\R$ and all $0<\rho_0<\rho_1<R$; and by
Step 2(ii), in the regime $\varepsilon<\lambda^2$ the first is available for
every $\ell\in[\mu/2,\mu]$ and the second for every $\ell\in[-\mu,-\mu/2]$.
Those are exactly the families Steps 6 and 7 consume.  Since
\noderef{DeGiorgiLevelBound} is weakened by increasing $\gamma$ --- its second
clause has $\gamma$ as a factor on the right --- a single $\gamma$ may serve all
the instances below, and we fix
\[
\gamma_\sharp=\max\bigl\{\gamma_9,\ 16\,C_{12}\bigr\} ,
\]
a function of $n,p,r$ alone, where $C_{12}$ is the constant of Step 4(4c); the
second entry is the one Step 6(6b) needs and is recorded here so that Step 5
can fix all its constants at once.

\emph{Step 4: the function $v$, its weak gradient, and inequality
\textup{(3.12)}.}  Assume for the whole of this step that
$\varepsilon<\lambda^2$, and set
\[
v=|g^k|^{p/2}\operatorname{sgn}(g^k),\qquad h_0=\lambda^{p/2},
\]
so that $|v|=|g^k|^{p/2}\le(2\lambda)^{p/2}=2^{p/2}h_0$ almost everywhere on
$B_R$ by (0b), which supplies its bound on the given $g$ in that form.  The
map $s\mapsto|s|^{p/2}\operatorname{sgn}s$ is a strictly increasing bijection of
$\R$ with inverse $t\mapsto|t|^{2/p}\operatorname{sgn}t$, so for every
$\ell\in\R$
\[
\{g^k>\ell\}=\{v>h\},\qquad
\{g^k<\ell\}=\{v<h\},\qquad
h=\operatorname{sgn}(\ell)|\ell|^{p/2} ,
\tag{4.1}
\]
and $|\ell|\le\lambda\iff|h|\le h_0$.

\textbf{(4a) $v\in W^{1,2}(B_T)$ for $p\ge2$.}  Let $F\colon\R\to\R$ agree with
$s\mapsto|s|^{p/2}\operatorname{sgn}s$ on $[-\mu,\mu]$ and be extended affinely
outside: $F(s)=\mu^{p/2}+\tfrac p2\mu^{p/2-1}(s-\mu)$ for $s>\mu$ and
$F(s)=-F(-s)$ for $s<-\mu$.  On $[-\mu,\mu]$ the derivative is
$\tfrac p2|s|^{p/2-1}$, which is continuous because $p/2-1\ge0$, and the two
affine pieces match $F$ and $F'$ at $\pm\mu$; hence $F\in C^1(\R)$ with
$|F'|\le L:=\tfrac p2\mu^{p/2-1}$ everywhere.  Fix $0<T<R$.  By (0.3),
$g^k\in W^{1,2}(B_T)$ with weak gradient $Dg^k$, and
$\operatorname{volume}B_T\ne\infty$, so \noderef{WeakGradientC1Comp} with
$q=2$, the bound $L$, the open set $B_T$ and this $F$ gives
\[
F\circ g^k\in W^{1,2}(B_T)\ \text{with weak gradient}\ x\mapsto F'(g^k(x))\,Dg^k(x).
\]
By (0b), $|g^k|\le\mu$ \emph{almost everywhere} on $B_R$ --- that is the form
(0b) supplies for the given $g$, the everywhere bound belonging to the
continuous representative $g'$ and not to $g$ --- and $B_T\subseteq B_R$, so
\[
F\circ g^k=v
\qquad\text{and}\qquad
F'(g^k)=\tfrac p2|g^k|^{p/2-1}
\qquad\text{almost everywhere on }B_T .
\tag{4.1a}
\]
The restriction to ``almost everywhere'' is not removable and is not a remark
about an irrelevant set: $F$ was built to agree with
$s\mapsto|s|^{p/2}\operatorname{sgn}s$ on $[-\mu,\mu]$ and to be affine outside
it, so at a point $x$ with $|g^k(x)|>\mu$ the two functions genuinely take
different values, and likewise $F'(g^k(x))=\tfrac p2\mu^{p/2-1}$ differs there
from $\tfrac p2|g^k(x)|^{p/2-1}$.  So (4.1a) is an almost-everywhere identity
between two \emph{pairs} of functions, and the membership has to be transported
across it rather than read off.

\emph{The transport, and both members move.}  By \noderef{MemW1pOn} at the
exponent $2$ on $B_T$, the assertion ``$f\in W^{1,2}(B_T)$ with weak gradient
$G$'' is the conjunction of three things: the defining identity of
\noderef{HasWeakGradientOn},
\[
\int_{B_T}f\,D_i\varphi\,dx=-\int_{B_T}G^i\,\varphi\,dx
\qquad(i=1,\dots,n),
\]
for every test function $\varphi$ of that node's class, together with
$f\in L^2(B_T)$ and $G\in L^2(B_T)$.  Each of the three is an integral
assertion: both sides of the identity are integrals whose integrands involve
$f$, respectively $G$, as a factor, and the two memberships are the finiteness
of $\int_{B_T}|f|^2dx$ and of $\int_{B_T}\|G\|^2dx$.  An integral is unchanged
when its integrand is altered on a null set, so each of the three holds for a
pair $(f,G)$ exactly when it holds for any pair agreeing with $(f,G)$ almost
everywhere on $B_T$.  This is the reading of those two definition nodes already
used at (0.3$'$) above.

Here \emph{both} members move, which is why the transport is stated for the
pair and not for the function alone: $v$ differs from $F\circ g^k$, and $V$
below differs from $F'(g^k)Dg^k$, on the same null set
$\{x\in B_T:|g^k(x)|>\mu\}$.  Applying the previous paragraph to the pair
$\bigl(F\circ g^k,\;F'(g^k)Dg^k\bigr)$, whose membership
\noderef{WeakGradientC1Comp} has just supplied, and to the pair $(v,V)$, which
agrees with it almost everywhere on $B_T$ by (4.1a), gives
$v\in W^{1,2}(B_T)$ with weak gradient
\[
V:=\tfrac p2|g^k|^{p/2-1}Dg^k,\qquad
\|V\|^2=\tfrac{p^2}4|g^k|^{p-2}\|Dg^k\|^2 ,
\tag{4.2}
\]
which is the paper's identity
$|\nabla(|D_ku|^{p/2})|^2=(p^2/4)|D_ku|^{p-2}|\nabla D_ku|^2$.  (At
$g^k(x)=0$ and $p=2$ the factor is $1$; for $p>2$ it is $0$.)

\textbf{(4b) $v\in W^{1,2}(B_T)$ for $1<p<2$, by the $\eta$-route.}  Here
$s\mapsto|s|^{p/2}\operatorname{sgn}s$ is not $C^1$ at $0$, and the paper's
regularisation is needed.  For $0<\eta$ put
\[
\pi_\eta(s)=(|s|+\eta)^{p-2}s,\qquad
G_\eta(s)=\operatorname{sgn}(s)\bigl[(|s|+\eta)^{p/2}-\eta^{p/2}\bigr],
\qquad
\psi_\eta=\pi_\eta\circ g^k .
\]
Both are odd and strictly increasing, and both are $C^1$ with bounded
derivative:
\[
\pi_\eta'(s)=(|s|+\eta)^{p-2}+(p-2)(|s|+\eta)^{p-3}|s|,
\qquad
G_\eta'(s)=\tfrac p2(|s|+\eta)^{p/2-1},
\]
the first continuous at $0$ because the second summand carries the factor $|s|$,
and both bounded, for $1<p<2$, by $\eta^{p-2}$ resp.\ $\tfrac p2\eta^{p/2-1}$,
since the exponents $p-2$ and $p/2-1$ are negative and $|s|+\eta\ge\eta$.  One
also has
\[
(p-1)(|s|+\eta)^{p-2}\ \le\ \pi_\eta'(s)\ \le\ (|s|+\eta)^{p-2} ,
\tag{4.3}
\]
the upper bound because $(p-2)(|s|+\eta)^{p-3}|s|\le0$, the lower because
$(|s|+\eta)^{p-3}|s|\le(|s|+\eta)^{p-2}$ and therefore
$\pi_\eta'(s)\ge(|s|+\eta)^{p-2}+(p-2)(|s|+\eta)^{p-2}=(p-1)(|s|+\eta)^{p-2}$.
These are the paper's inequalities on p.~21.  By (0.3) and
\noderef{WeakGradientC1Comp} on $B_T$, exactly as in (4a),
\[
\psi_\eta\in W^{1,2}(B_T)\ \text{with weak gradient}\ \pi_\eta'(g^k)Dg^k,
\qquad
v_\eta:=G_\eta\circ g^k\in W^{1,2}(B_T)\ \text{with weak gradient}
\]
\[
V_\eta:=\tfrac p2(|g^k|+\eta)^{p/2-1}Dg^k,
\qquad
\|V_\eta\|^2=\tfrac{p^2}4(|g^k|+\eta)^{p-2}\|Dg^k\|^2 ,
\tag{4.4}
\]
which is (3.14) of the thesis.  Note $|\psi_\eta|\le(|g^k|+\eta)^{p-2}|g^k|\le
|g^k|^{p-1}\le(2\lambda)^{p-1}$, using $p-2<0$, and $|v_\eta|\le(2\lambda+\eta)^{p/2}$.

\emph{The differentiated equation against the regularised truncation.}  Fix
$\ell\in\R$, $0<\sigma<1$ and $\rho\le R$, and let
$\zeta=\zeta_{\rho-\sigma\rho,\rho}$ of (0e).  Choose $T<R$ with
$\operatorname{tsupport}\zeta\subseteq B_T$ (possible since
$\operatorname{tsupport}\zeta$ is a compact subset of $B_\rho\subseteq B_R$).
Put
\[
\theta^+_\eta=(\psi_\eta-\ell)^+,\qquad \varphi^+=\theta^+_\eta\zeta^2 .
\]
By \noderef{WeakGradientPositivePart} with $q=2$, $c=\ell$, the open set $B_T$
of finite volume and the pair $(\psi_\eta,\pi_\eta'(g^k)Dg^k)$,
$\theta^+_\eta\in W^{1,2}(B_T)$ with weak gradient
$\chi_{\{\psi_\eta>\ell\}}\pi_\eta'(g^k)Dg^k$, which is the paper's
$(\theta^+_\eta)'=\pi_\eta'\chi_{\{\pi_\eta>\ell\}}$.  Since this holds for every
$T$ with $\operatorname{tsupport}\zeta\subseteq B_T<R$ and in fact for every
$T<R$ (the construction does not use $\zeta$), (0.3$'$) gives
$\theta^+_\eta\in W^{1,2}_{\mathrm{loc}}(B_R)$ with that weak gradient.  Now
$\zeta^2$ is smooth with compact support in $B_R$, so clause (2) of
\noderef{WeakGradientSmoothMul} with $q=2$ and $\Omega=B_R$ gives
$\varphi^+\in W^{1,2}(B_R)$ with weak gradient
\[
g_{\varphi^+}=\zeta^2\chi_{\{\psi_\eta>\ell\}}\pi_\eta'(g^k)Dg^k
+\theta^+_\eta\,\nabla(\zeta^2),
\qquad \nabla(\zeta^2)=2\zeta\nabla\zeta ,
\]
and $\varphi^+$ has compact support with
$\operatorname{tsupport}\varphi^+\subseteq\operatorname{tsupport}\zeta\subseteq
B_R$.  Apply \noderef{DifferentiatedEquationWeak} with
$n,p,\varepsilon,\gamma_0,\gamma_1$ --- its hypotheses $2\le n$, $1<p$,
$0<\varepsilon$, $0<\gamma_0$ and $\gamma_0\le\gamma_1$, the last by (0c) ---
the open set $B_R$, the field $A$, the solution $u$ with gradient $g$, the
family $Dg$, the index $k$ and the pair $(\varphi^+,g_{\varphi^+})$.  It gives
\[
\int_{B_R}\sum_{i,j}\bigl(DA(g)e_j\bigr)_i\,(Dg^k)_j\,(g_{\varphi^+})_i\,dx=0 ,
\]
i.e., separating the two summands of $g_{\varphi^+}$,
\[
\int_{B_R}\chi_{\{\psi_\eta>\ell\}}\pi_\eta'(g^k)\,
\bigl\langle DA(g)Dg^k,Dg^k\bigr\rangle\zeta^2
=-2\int_{B_R}\theta^+_\eta\,\zeta\,\bigl\langle DA(g)Dg^k,\nabla\zeta\bigr\rangle .
\]
On the left, \texttt{Q1} of \noderef{EllipticStructureQ12} with $\xi=Dg^k$ and
$h=g(x)$ --- its guard $\varepsilon+\|g\|^2=w\ne0$ holds since
$\varepsilon>0$ --- gives $\langle DA(g)Dg^k,Dg^k\rangle\ge\gamma_0
w^{p/2-1}\|Dg^k\|^2$, and (4.3) gives
$\pi_\eta'(g^k)\ge(p-1)(|g^k|+\eta)^{p-2}$.  On the right, \texttt{Q2} gives
$|(DA(g)e_j)_i|\le\gamma_1w^{p/2-1}$, whence
\[
\bigl|\langle DA(g)Dg^k,\nabla\zeta\rangle\bigr|
\le\gamma_1w^{\frac p2-1}\sum_{i,j}|(Dg^k)_j|\,|(\nabla\zeta)_i|
\le n\gamma_1w^{\frac p2-1}\|Dg^k\|\,|\nabla\zeta| ,
\]
using $\sum_j|a_j|\le\sqrt n\,\|a\|$ (Cauchy--Schwarz, Preamble) twice.  Hence
\[
\gamma_0(p-1)\!\!\int_{A^+_{\ell}(\eta)}\!\!w^{\frac p2-1}\|Dg^k\|^2(|g^k|+\eta)^{p-2}\zeta^2
\le2n\gamma_1\!\!\int_{B_R}\!\!w^{\frac p2-1}\|Dg^k\|\,\theta^+_\eta\,|\zeta|\,|\nabla\zeta| ,
\tag{3.13$^+$}
\]
where $A^+_{\ell}(\eta)=B_R\cap\{\psi_\eta>\ell\}$; the integrand on the right
vanishes off $A^+_\ell(\eta)$ because $\theta^+_\eta$ does.  This is
(3.13) of the thesis.

\emph{The Cauchy--Schwarz step.}  Assume $\eta^2<\varepsilon$.  Then
$w\ge\varepsilon>\eta^2$ and $w\ge(g^k)^2$, so
$w^{1/2}\ge\max\{\eta,|g^k|\}\ge\tfrac12(|g^k|+\eta)$ and therefore, since
$1-p/2>0$,
\[
(|g^k|+\eta)^{1-\frac p2}\le\bigl(2w^{1/2}\bigr)^{1-\frac p2}
=2^{1-\frac p2}w^{\frac{2-p}4}.
\]
From (4.4), $\|Dg^k\|=\tfrac2p(|g^k|+\eta)^{1-\frac p2}\|V_\eta\|$, so
\[
w^{\frac p2-1}\|Dg^k\|\le\tfrac2p2^{1-\frac p2}\,
w^{\frac p2-1+\frac{2-p}4}\|V_\eta\|
=C(p)\,w^{\frac{p-2}4}\|V_\eta\| ,
\tag{3.15}
\]
because $\tfrac p2-1+\tfrac{2-p}4=\tfrac{p-2}4$.  Substituting (3.15) and
(4.4) into (3.13$^+$) and applying the Cauchy--Schwarz inequality (Preamble) to
the right-hand side, written as the product of
$w^{\frac{p-2}4}\|V_\eta\|\,|\zeta|$ and $\theta^+_\eta|\nabla\zeta|$ and noting
$\bigl(w^{\frac{p-2}4}\bigr)^2=w^{\frac p2-1}$,
\[
\tfrac{4\gamma_0(p-1)}{p^2}\!\!\int_{A^+_\ell(\eta)}\!\!w^{\frac p2-1}\|V_\eta\|^2\zeta^2
\le2n\gamma_1C(p)\Bigl(\int_{A^+_\ell(\eta)}\!\!w^{\frac p2-1}\|V_\eta\|^2\zeta^2\Bigr)^{\frac12}
\Bigl(\int_{B_R}(\theta^+_\eta)^2|\nabla\zeta|^2\Bigr)^{\frac12}.
\]
The first factor on the left is finite --- $\|V_\eta\|\le\tfrac p2\eta^{p/2-1}\|Dg^k\|$
is in $L^2$ of the compact $\operatorname{tsupport}\zeta$ by
\noderef{MemW1pLocOn}, and $w^{p/2-1}\le\varepsilon^{p/2-1}$ since $p/2-1<0$ and
$w\ge\varepsilon$ --- so it may be divided out, and squaring gives, with
$\gamma_1\le r\gamma_0$,
\[
\int_{A^+_\ell(\eta)}w^{\frac p2-1}\|V_\eta\|^2\zeta^2\,dx
\ \le\ C_{16}\int_{B_R}\bigl((\psi_\eta-\ell)^+\bigr)^2|\nabla\zeta|^2\,dx ,
\tag{3.16$^+$}
\]
with $C_{16}=\bigl(2nrC(p)p^2/(4(p-1))\bigr)^2$, a function of $n,p,r$.  The
$-$ case (3.16$^-$), with $(\ell-\psi_\eta)^+$ and
$A^-_\ell(\eta)=B_R\cap\{\psi_\eta<\ell\}$, follows by the same argument from
$\varphi^-=(\ell-\psi_\eta)^+\zeta^2$, using
\noderef{WeakGradientPositivePart} applied to $(-\psi_\eta)$ at the level
$-\ell$ (so that $\theta^-_\eta$ has weak gradient
$-\chi_{\{\psi_\eta<\ell\}}\pi_\eta'(g^k)Dg^k$) together with
\noderef{WeakGradientLinear}.

\emph{The limit $\eta\downarrow0$.}  Put
\[
\psi=|g^k|^{p-2}g^k\ \text{on}\ \{g^k\ne0\},\quad\psi=0\ \text{on}\ \{g^k=0\},
\qquad
V=\tfrac p2|g^k|^{p/2-1}Dg^k\ \text{on}\ \{g^k\ne0\},\quad V=0\ \text{on}\ \{g^k=0\} .
\]
We claim $Dg^k=0$ a.e.\ on $\{g^k=0\}$.  Indeed on $B_T$,
\noderef{WeakGradientPositivePart} at the level $c=0$ gives
$(g^k)^+\in W^{1,2}(B_T)$ with weak gradient $\chi_{\{g^k>0\}}Dg^k$, and
applied to $-g^k$ it gives $(-g^k)^+$ with weak gradient
$-\chi_{\{g^k<0\}}Dg^k$; since $g^k=(g^k)^+-(-g^k)^+$,
\noderef{WeakGradientLinear} (with $a=1$, $b=-1$) shows that
$\chi_{\{g^k>0\}}Dg^k+\chi_{\{g^k<0\}}Dg^k$ is also a weak gradient of $g^k$ on
$B_T$; and a weak gradient is determined almost everywhere (Preamble: the
fundamental lemma of the calculus of variations, second consequence).  Hence
$\chi_{\{g^k=0\}}Dg^k=0$ a.e., as claimed; so
$\|V_\eta\|\uparrow\|V\|$ pointwise a.e.\ as $\eta\downarrow0$, the exponent
$p/2-1$ being negative on $\{g^k\ne0\}$ and both sides being $0$ a.e.\ on
$\{g^k=0\}$.

\emph{First, $v\in W^{1,2}(B_T)$ with weak gradient $V$.}  Take $\ell=0$ in
(3.16$^\pm$): then $A^+_0(\eta)=B_R\cap\{g^k>0\}$ and
$A^-_0(\eta)=B_R\cap\{g^k<0\}$ for every $\eta$, and
$(\psi_\eta^{\pm})^2\le(2\lambda)^{2(p-1)}$.  Fix $T<R$ and
$\zeta=\zeta_{T,T'}$ with $T<T'<R$; then $\zeta=1$ on $\overline B_T$ and
$|\nabla\zeta|\le\Lambda/(T'-T)$, so (3.16$^\pm$) and
$w^{p/2-1}\ge\bigl((1+4n)\lambda^2\bigr)^{p/2-1}>0$ (by (0.1), $p/2-1<0$ and
$\varepsilon<\lambda^2$) give
\[
\int_{B_T\cap\{g^k\ne0\}}\|V_\eta\|^2
\ \le\ \bigl((1+4n)\lambda^2\bigr)^{1-\frac p2}C_{16}
\frac{\Lambda^2}{(T'-T)^2}(2\lambda)^{2(p-1)}\,2|B_R|\ =:\ \mathcal K<\infty
\]
for every $\eta$ with $\eta^2<\varepsilon$.  By Fatou's lemma (Preamble) and
$\|V_\eta\|\uparrow\|V\|$, $\int_{B_T}\|V\|^2\le\mathcal K<\infty$.  Next, restricting to $\eta\le1$,
$v_\eta\to v$ pointwise with $|v_\eta|\le(2\lambda+1)^{p/2}$ and
$V_\eta\to V$ pointwise a.e.\ with $\|V_\eta\|\le\|V\|\in L^2(B_T)$, so for
every test function $\varphi$ with $\operatorname{tsupport}\varphi\subseteq B_T$
dominated convergence (Preamble) passes to the limit in the identity of
\noderef{HasWeakGradientOn} for $(v_\eta,V_\eta)$ and yields it for $(v,V)$;
together with $|v|\le2^{p/2}h_0$ and $\|V\|\in L^2(B_T)$ this is
$v\in W^{1,2}(B_T)$ with weak gradient $V$ (\noderef{MemW1pOn}), and (4.2)
holds verbatim.

\emph{Second, (3.17) of the thesis, for every $\ell\in\R$.}  Write
$A^+_\ell=B_R\cap\{\psi>\ell\}$ and $A^+_\ell(\eta)=B_R\cap\{\psi_\eta>\ell\}$.
Fatou's lemma (Preamble) needs only a pointwise lower bound on the liminf of the
integrands, and \emph{not} any monotonicity of the sets $A^+_\ell(\eta)$; what
has to be checked is therefore the pointwise statement
\[
x\in A^+_\ell\ \Longrightarrow\ x\in A^+_\ell(\eta)\ \text{ for all sufficiently
small }\eta>0 ,
\tag{3.16$'$}
\]
and that holds for \emph{every} real $\ell$.  Indeed let $x\in B_R$ with
$\psi(x)>\ell$ and put $t=g^k(x)$, so that $\psi(x)=|t|^{p-2}t$ and
$\psi_\eta(x)=(|t|+\eta)^{p-2}t$; recall $p-2<0$.  Three cases, and only the
first needs a limit at all.
\begin{itemize}
\item $t>0$.  Then $\eta\mapsto(t+\eta)^{p-2}t$ is continuous at $\eta=0$,
  because $t>0$ keeps the base of the power away from $0$, and its value there is
  $|t|^{p-2}t=\psi(x)>\ell$; so $\psi_\eta(x)>\ell$ for all small $\eta>0$.
\item $t=0$.  Then $\psi(x)=|0|^{p-2}\cdot0=0$ and
  $\psi_\eta(x)=\eta^{p-2}\cdot0=0$ for every $\eta$ --- the convention
  $0^{p-2}=0$ for $p\ne2$ is the run's rpow convention and is irrelevant here,
  since the factor $t$ vanishes.  So $\psi(x)>\ell$ says $\ell<0$, and
  $\psi_\eta(x)=0>\ell$ for \emph{every} $\eta>0$.
\item $t<0$.  Then $|t|+\eta>|t|>0$ and $p-2<0$, so
  $(|t|+\eta)^{p-2}\le|t|^{p-2}$; multiplying by $t<0$ reverses the inequality
  and gives $\psi_\eta(x)\ge\psi(x)>\ell$ for \emph{every} $\eta>0$.
\end{itemize}
It is worth recording what is \emph{not} true, because it is what a
monotonicity-based argument would need: on $\{g^k>0\}$ one has
$\psi_\eta\le\psi$, by the computation of the third bullet run at $t>0$, where
multiplying $(|t|+\eta)^{p-2}\le|t|^{p-2}$ by $t>0$ preserves the direction; so the
regularisation sits strictly \emph{below} its limit there and the sets
$A^+_\ell(\eta)$ do not increase to $A^+_\ell$.  That costs nothing, because
Fatou consumes (3.16$'$) and not inclusion.

Now apply Fatou's lemma to the nonnegative integrands
$\chi_{A^+_\ell(\eta)}w^{p/2-1}\|V_\eta\|^2\zeta^2$ along any sequence
$\eta_j\downarrow0$: at each $x\in A^+_\ell$ the characteristic factor equals $1$
for all large $j$ by (3.16$'$) and $\|V_{\eta_j}(x)\|\uparrow\|V(x)\|$ as shown
above, so the liminf of the integrands is at least
$\chi_{A^+_\ell}w^{p/2-1}\|V\|^2\zeta^2$ at every $x$; hence the integral of the
latter is at most the liminf of (3.16$^+$)'s left-hand sides.  Meanwhile
$(\psi_\eta-\ell)^+\to(\psi-\ell)^+$ pointwise everywhere: at a point with
$t\ne0$ the map $\eta\mapsto(|t|+\eta)^{p-2}t$ is continuous at $\eta=0$ with
value $\psi(x)$, since $|t|>0$, and at a point with $t=0$ one has
$\psi_\eta(x)=0=\psi(x)$ for every $\eta$; so $\psi_\eta\to\psi$ pointwise, and
$s\mapsto(s-\ell)^+$ is continuous.  (This pointwise convergence alone already
gives (3.16$'$) for every $\ell$; the three cases above are kept because they
show which of them needs a limit and which does not, and because the case $t=0$
is the one where the rpow convention has to be read.)  The bound is
$\bigl((\psi_\eta-\ell)^+\bigr)^2\le\bigl((2\lambda)^{p-1}+|\ell|\bigr)^2$,
since $|\psi_\eta|\le(|g^k|+\eta)^{p-2}|g^k|\le|g^k|^{p-1}\le(2\lambda)^{p-1}$
almost everywhere, by $p-2<0$ and (0b) in the almost-everywhere form it supplies
for $g$ --- which is the form dominated convergence asks for, its domination
hypothesis being almost everywhere --- and $|\nabla\zeta|^2$ is integrable on
$B_R$; so dominated
convergence (Preamble) identifies the limit of the right-hand sides.  Therefore
\[
\int_{A^+_\ell}w^{\frac p2-1}\|V\|^2\zeta^2\,dx
\ \le\ C_{16}\int_{B_R}\bigl((\psi-\ell)^+\bigr)^2|\nabla\zeta|^2\,dx
\qquad(\ell\in\R),
\tag{3.17$^+$}
\]
and symmetrically (3.17$^-$) for every $\ell\in\R$, with $A^-_\ell=B_R\cap
\{\psi<\ell\}$ and $(\ell-\psi)^+$ in place of $(\psi-\ell)^+$: there the three
cases are $t<0$ (continuity, $|t|>0$), $t=0$ (then $\psi(x)=0<\ell$ forces
$\ell>0$ and $\psi_\eta(x)=0<\ell$ for every $\eta$) and $t>0$ (where
$\psi_\eta\le\psi<\ell$ for every $\eta$).  The full range $\ell\in\R$ is the
paper's own: (3.17) of the thesis at thesis l.~779--781 carries no restriction and
(3.18) of the thesis is asserted ``for all $\ell\in\R$'' at l.~793.

\textbf{(4c) Inequality (3.12) of the thesis.}  We prove, with the ranges the paper
prints at thesis l.~728--729: for every $0<\sigma<1$, every $\rho\le R$ and
every $h\ge-h_0$,
\[
\int_{B_{\rho-\sigma\rho}\cap\{v>h\}}\|V\|^2\,dx
\ \le\ C_{12}\,\frac{h_0^2}{(\sigma\rho)^2}\bigl|B_\rho\cap\{v>h\}\bigr| ,
\tag{3.12$^+$}
\]
and, for every $h\le h_0$, the same bound with $\{v<h\}$ in place of
$\{v>h\}$ on both sides,
\[
\int_{B_{\rho-\sigma\rho}\cap\{v<h\}}\|V\|^2\,dx
\ \le\ C_{12}\,\frac{h_0^2}{(\sigma\rho)^2}\bigl|B_\rho\cap\{v<h\}\bigr| .
\tag{3.12$^-$}
\]
Here $C_{12}=C_{12}(n,p,r)$.  Let $\zeta=\zeta_{\rho-\sigma\rho,\rho}$ of (0e),
so $|\nabla\zeta|\le\Lambda/(\sigma\rho)$, $\zeta=1$ on
$\overline B_{\rho-\sigma\rho}$ and $\operatorname{supp}\zeta\subseteq B_\rho$.
In both regimes the level $h$ is transported to a level $\ell$ for $g^k$ resp.\
for $\psi$ by the strictly increasing bijection that carries one to the other,
and \emph{only a one-sided bound on $\ell$ is used}: the $+$ case needs a lower
bound for $\ell$ and the $-$ case an upper one.  That is why the two ranges are
one-sided, and it is what the paper prints.

\emph{Case $p\ge2$, the $+$ side.}  Let $h\ge-h_0$ and put
$\ell=\operatorname{sgn}(h)|h|^{2/p}$, so that $\{v>h\}=\{g^k>\ell\}$ by (4.1)
and $\ell\ge-\lambda$: for $h\ge0$ this is clear since $\ell\ge0$, and for
$-h_0\le h<0$ it is $|\ell|=|h|^{2/p}\le h_0^{2/p}=\lambda$.  In (3.4$^+$) bound
the left-hand side below using $w\ge(g^k)^2$ and $p/2-1\ge0$, so
$w^{p/2-1}\ge|g^k|^{p-2}$, and then (4.2):
\[
\int_{B_R\cap\{g^k>\ell\}}w^{\frac p2-1}\|Dg^k\|^2\zeta^2
\ \ge\ \int_{B_R\cap\{g^k>\ell\}}|g^k|^{p-2}\|Dg^k\|^2\zeta^2
\ =\ \frac4{p^2}\int_{B_R\cap\{v>h\}}\|V\|^2\zeta^2 .
\]
On the right-hand side of (3.4$^+$), the integrand is supported in
$\{g^k>\ell\}\cap B_\rho$; there $w\le(1+4n)\lambda^2$ by (0.1) and
$\varepsilon<\lambda^2$, so $w^{p/2-1}\le(1+4n)^{p/2-1}\lambda^{p-2}$ (again
$p/2-1\ge0$), and $0<g^k-\ell\le\mu-\ell\le2\lambda+\lambda=3\lambda$ because
$\ell\ge-\lambda$.  Hence
\[
c_4\int_{B_R}w^{\frac p2-1}\bigl((g^k-\ell)^+\bigr)^2|\nabla\zeta|^2
\le c_4(1+4n)^{\frac p2-1}\lambda^{p-2}\,9\lambda^2\,
\frac{\Lambda^2}{(\sigma\rho)^2}\bigl|B_\rho\cap\{v>h\}\bigr| .
\]
Since $\lambda^{p-2}\lambda^2=\lambda^p=(\lambda^{p/2})^2=h_0^2$, and since
$\zeta=1$ on $\overline B_{\rho-\sigma\rho}$, this is (3.12$^+$) with
$C_{12}=\tfrac{9p^2}4c_4(1+4n)^{p/2-1}\Lambda^2$.

\emph{Case $p\ge2$, the $-$ side.}  Let $h\le h_0$ and put the same
$\ell=\operatorname{sgn}(h)|h|^{2/p}$, so that $\{v<h\}=\{g^k<\ell\}$ and now
$\ell\le\lambda$, by the mirror image of the computation just made.  Run the
same two estimates on (3.4$^-$), whose right-hand integrand is supported in
$\{g^k<\ell\}\cap B_\rho$ and satisfies
$0<\ell-g^k\le\ell+\mu\le\lambda+2\lambda=3\lambda$; this is (3.12$^-$) with the
same $C_{12}$.

\emph{Case $1<p<2$, the $-$ side.}  Let $h\le h_0$ and put
\[
\ell=\operatorname{sgn}(h)|h|^{\frac{2(p-1)}p},
\qquad\text{so}\qquad
h=\operatorname{sgn}(\ell)|\ell|^{\frac p{2(p-1)}} .
\]
Then $\{\psi<\ell\}=\{g^k<\operatorname{sgn}(\ell)|\ell|^{1/(p-1)}\}=\{v<h\}$,
because $s\mapsto|s|^{p-2}s$ and $s\mapsto|s|^{p/2}\operatorname{sgn}s$ are
strictly increasing bijections of $\R$ and the two displayed exponents are
inverse to one another on the relevant branch; and the level transports exactly:
\[
h\le h_0=\lambda^{p/2}
\iff
\ell\le\lambda^{p-1},
\qquad\text{because}\qquad
\bigl(\lambda^{p-1}\bigr)^{\frac p{2(p-1)}}=\lambda^{\frac p2}
\]
and $t\mapsto\operatorname{sgn}(t)|t|^{p/(2(p-1))}$ is increasing.  This is the
identity that makes the $\eta$-route land on (3.12) rather than on a weaker
relative of it; the route's bookkeeping is in $\ell$ and (3.12)'s is in $h$, and
$\lambda^{p-1}$ must not be carried into an $h$-indexed statement.
Now use (3.17$^-$) at this $\ell$ --- available for \emph{every} real $\ell$ by
the preceding paragraph, which is the point of its repair.  Its left-hand side:
$w^{p/2-1}\ge\bigl((1+4n)\lambda^2\bigr)^{p/2-1}=(1+4n)^{p/2-1}\lambda^{p-2}$ by
(0.1), $\varepsilon<\lambda^2$ and $p/2-1<0$, and $\zeta=1$ on
$\overline B_{\rho-\sigma\rho}$, so
\[
(1+4n)^{\frac p2-1}\lambda^{p-2}\int_{B_{\rho-\sigma\rho}\cap\{v<h\}}\|V\|^2
\ \le\ \int_{A^-_\ell}w^{\frac p2-1}\|V\|^2\zeta^2 ,
\]
which is the left-hand side of (3.19) of the thesis.  Its right-hand side: the
integrand is supported in $\{\psi<\ell\}\cap B_\rho$, where
$0<\ell-\psi\le\ell^++(2\lambda)^{p-1}\le\lambda^{p-1}+2^{p-1}\lambda^{p-1}
=(1+2^{p-1})\lambda^{p-1}$, using $\ell\le\lambda^{p-1}$ and
$|\psi|\le(2\lambda)^{p-1}$, so
\[
C_{16}\int_{B_R}\bigl((\ell-\psi)^+\bigr)^2|\nabla\zeta|^2
\le C_{16}(1+2^{p-1})^2\lambda^{2p-2}\frac{\Lambda^2}{(\sigma\rho)^2}
\bigl|B_\rho\cap\{v<h\}\bigr| .
\]
Dividing both sides by $(1+4n)^{p/2-1}\lambda^{p-2}>0$ --- strictly positive for
every real $p$ because $\lambda>0$, including the present case where the
exponent $p-2$ is negative --- and using
$\lambda^{2p-2}/\lambda^{p-2}=\lambda^{p}=h_0^2$, we obtain (3.12$^-$) with
$C_{12}=C_{16}(1+2^{p-1})^2(1+4n)^{1-p/2}\Lambda^2$.

\emph{Case $1<p<2$, the $+$ side.}  Let $h\ge-h_0$ and put the same $\ell$, so
that $\{\psi>\ell\}=\{v>h\}$ and $\ell\ge-\lambda^{p-1}$ by the mirror image of
the transport identity.  Run the same computation from (3.17$^+$) at this
$\ell$ --- again available for every real $\ell$ --- the right-hand bound now
being $0<\psi-\ell\le(2\lambda)^{p-1}+\ell^-\le(1+2^{p-1})\lambda^{p-1}$ with
$\ell^-=\max\{-\ell,0\}\le\lambda^{p-1}$.  This gives (3.12$^+$).

Enlarging $C_{12}$ to the maximum of the two expressions above, (3.12$^\pm$)
holds in both regimes, on the full one-sided ranges $h\ge-h_0$ and $h\le h_0$,
with one constant depending only on $n,p,r$.

\emph{Step 5: the constants $\theta_2$ and $\nu$.}  All of the radii below are
$<R$, because \noderef{LevelTruncationMaximumPrinciple} and
\noderef{LevelSetMeasureImprovement} ask for $W^{1,2}$ membership on the
\emph{whole} of their outer ball (\noderef{MemW1pOn}), whereas the present
lemma assumes only $W^{1,2}_{\mathrm{loc}}(B_R)$, which is what the cited
regularity (C1) supplies; (0.3) converts the latter into the former on every
$B_T$ with $T<R$, and the radii are chosen so that the hypotheses still follow
from those stated at radius $R$ and the conclusions still land on $B_{R/2}$.
Put
\[
S=\tfrac34R,\qquad \lambda_1'=\frac R{2S}=\frac23\in(0,1),
\qquad\text{so}\qquad \lambda_1'S=\tfrac R2 ,
\]
and $|B_S|=(3/4)^n|B_R|$ (Preamble: the volume of a ball).  Set
\[
\lambda_0^{(1)}=\tfrac34,\qquad
\lambda_0^{(2)}=\frac{2^{p/2}+\tfrac12}{2^{p/2}+1},
\]
both in $(0,1)$.  Apply \noderef{LevelTruncationMaximumPrinciple} twice, with
$n$, the constant $\gamma_\sharp$ of Step 3, $\lambda_0=\lambda_0^{(1)}$ resp.\
$\lambda_0^{(2)}$, and $\lambda_1=\lambda_1'$, its hypothesis
\noderef{CitedSobolevPoincareBall} being the present lemma's; let $\theta^{(1)}$
and $\theta^{(2)}$ be the constants $\theta_0$ it produces, and put
\[
\theta_0^\ast=\min\{\theta^{(1)},\theta^{(2)},1\}>0,
\qquad
\theta_2=1-\theta_0^\ast(3/4)^n\ \in(0,1) .
\]
(Decreasing $\theta_0$ only strengthens the hypothesis of that node, so
$\theta_0^\ast$ may be used in place of either constant.)  Next put
\[
t=\Bigl(\frac{1+\theta_2}2\Bigr)^{1/n}\in(\theta_2^{1/n},1),
\qquad
S_b=\frac{1+t}2R,\qquad
\lambda_1=\frac{2t}{1+t}\in(0,1),
\qquad\text{so}\qquad \lambda_1S_b=tR,
\]
and note $t>2^{-1/n}>\tfrac12$ for $n\ge2$, so $S_b\in(R/2,R)$ and
$\lambda_1S_b>R/2$; also $\lambda_1<1$ because $t<1$.  Put
\[
c_0=\frac{\theta_2}{t^n}=\frac{2\theta_2}{1+\theta_2}\in[0,1),
\qquad
\lambda_1''=\frac R{2\lambda_1S_b}=\frac1{2t}\in(\tfrac12,1) .
\]
Apply \noderef{LevelTruncationMaximumPrinciple} once more, with $n$,
$\gamma_\sharp$, $\lambda_0=\tfrac12$ and $\lambda_1=\lambda_1''$, and let
$\theta^{(3)}$ be its constant; put $\theta_1=\min\{\theta^{(3)},\tfrac12\}\in(0,1)$.
Apply \noderef{LevelSetMeasureImprovement} with $n$, $\gamma_\sharp$,
$\lambda_1$, $c_0$ and $\theta_1$, and let $\lambda_2\in(0,1)$ be the constant
it produces.  Finally put
\[
\nu=1-\frac{\lambda_2}4\ \in(0,1).
\]
Every one of $S,\lambda_1',\lambda_0^{(1)},\lambda_0^{(2)},\theta_0^\ast,
\theta_2,t,S_b/R,\lambda_1,c_0,\lambda_1'',\theta_1,\lambda_2,\nu$ is a function
of $n$, $p$ and $\gamma_\sharp$, hence of $n,p,r$ alone, and all are chosen
before $\varepsilon,\gamma_0,\gamma_1,\mu,A,u,g,Dg,x_0,R$, which is the
quantifier order the statement asserts.  The order of choice is acyclic:
$\theta^{(1)},\theta^{(2)}$ depend on $\lambda_1'=2/3$, which is fixed in
advance; $\theta_2$ on those; $t,\lambda_1,c_0,\lambda_1''$ on $\theta_2$; and
$\theta^{(3)},\theta_1,\lambda_2,\nu$ on those in turn.

\emph{Step 6: part \textup{(a)}.}  Fix $k$ and assume the hypothesis of (a).
Since $\{g^k>\mu/2\}$ and $\{g^k<-\mu/2\}$ are disjoint, in the first
alternative
\[
\bigl|B_S\cap\{g^k<\tfrac\mu2\}\bigr|
\le\bigl|B_R\cap\{g^k<\tfrac\mu2\}\bigr|
\le|B_R|-\bigl|B_R\cap\{g^k>\tfrac\mu2\}\bigr|
\le(1-\theta_2)|B_R|=\theta_0^\ast|B_S| ,
\tag{6.1}
\]
and in the second, with the two sets exchanged,
$|B_S\cap\{g^k>-\mu/2\}|\le\theta_0^\ast|B_S|$.

\textbf{(6a) The case $\varepsilon\ge\lambda^2$.}  Treat the first alternative;
the second is the same argument with $g^k$ in place of $-g^k$.  Apply
\noderef{LevelTruncationMaximumPrinciple} --- the instance of Step 5 at
$\lambda_0^{(1)}=\tfrac34$ and $\lambda_1'=\tfrac23$ --- with the centre $x_0$,
the radius $S$, the level $\kappa=-\mu/2$, the function $f=-g^k$ and the field
$-Dg^k$, and $M=\operatorname{ess\,sup}_{B_S}(-g^k)$.  Its hypotheses: $0<S$;
$-g^k\in W^{1,2}(B_S)$ with weak gradient $-Dg^k$, by (0.3) and
\noderef{WeakGradientLinear}; $-g^k\le\mu$ a.e.\ on $B_S$, which is the required
essential upper bound; $M$ is by definition the essential supremum; the measure
hypothesis is (6.1), since
$\{\kappa<-g^k\}=\{g^k<\mu/2\}$ and $\theta_0^\ast\le\theta^{(1)}$; and the
family of $(2.1)$ instances is required for the levels
$l\in[\kappa,\tfrac34\kappa+\tfrac14M]$ and the radii
$\rho_0,\rho_1\in[\lambda_1'S,S]=[R/2,3R/4]$ with $\rho_0<\rho_1$, which Step 3
supplies from Step 2(i) --- available for \emph{all} levels in this case --- with
$\gamma=\gamma_\sharp$ and $\rho_1\le 3R/4<R$.  If $M\le\kappa$ then
$-g^k\le-\mu/2$ a.e.\ on $B_S\supseteq B_{R/2}$, i.e.\ $|g^k|\ge\mu/2$ a.e.\
there, and we are done; so assume $\kappa<M$, as the node requires.  Its
conclusion is
\[
-g^k\le\tfrac34\bigl(-\tfrac\mu2\bigr)+\tfrac14M
\le-\tfrac{3\mu}8+\tfrac\mu4=-\tfrac\mu8
\qquad\text{a.e.\ on }B_{\lambda_1'S}=B_{R/2},
\]
using $M\le\mu$.  Hence $g^k\ge\mu/8$ and so
$|g^k|\ge\tfrac18\mu\ge\min\{\tfrac1{16},2^{-1-4/p}\}\mu$ a.e.\ on $B_{R/2}$,
because $\min\{1/16,2^{-1-4/p}\}\le1/16\le1/8$.

\textbf{(6b) The case $\varepsilon<\lambda^2$.}  Treat the second alternative,
$|B_R\cap\{g^k<-\mu/2\}|\ge\theta_2|B_R|$; the first is the same argument
applied to $-g^k$, whose associated function is $-v$.  Let $v$, $V$, $h_0$ be as
in Step 4 and let
$M_k=\operatorname{ess\,sup}_{B_{R/2}}v$.

\emph{Branch 1: $M_k<-h_0/4$.}  Then $v<-h_0/4$ a.e.\ on $B_{R/2}$, i.e.\
$|g^k|^{p/2}\operatorname{sgn}(g^k)<-\tfrac14\lambda^{p/2}$; in particular
$g^k<0$ and $|g^k|^{p/2}>\tfrac14\lambda^{p/2}$, so
$|g^k|>4^{-2/p}\lambda=2^{-1-4/p}\mu$ a.e.\ on $B_{R/2}$, which is the
conclusion of (a) --- and this is precisely where the constant
$2^{-1-4/p}$ arises; see the Correction below.

\emph{Branch 2: $M_k\ge-h_0/4$.}  By (4.1),
$\{v<-h_0\}=\{g^k<-\lambda\}=\{g^k<-\mu/2\}$, so the second alternative reads
$|B_R\cap\{v<-h_0\}|\ge\theta_2|B_R|$ and, exactly as in (6.1),
\[
\bigl|B_S\cap\{v>-h_0\}\bigr|\le\theta_0^\ast|B_S| .
\tag{6.2}
\]
Apply \noderef{LevelTruncationMaximumPrinciple} --- the instance of Step 5 at
$\lambda_0^{(2)}$ and $\lambda_1'=\tfrac23$ --- with the centre $x_0$, the
radius $S$, the level $\kappa=-h_0$, the function $f=v$, the field $V$ and
$M=\operatorname{ess\,sup}_{B_S}v$.  Its hypotheses: $v\in W^{1,2}(B_S)$ with
weak gradient $V$, by Step 4 ((4a) for $p\ge2$, (4b) for $1<p<2$) at $T=S<R$;
$v\le2^{p/2}h_0$ a.e.; the measure hypothesis (6.2); and the family of $(2.1)$
instances for the levels
$l\in[-h_0,\lambda_0^{(2)}(-h_0)+(1-\lambda_0^{(2)})M]$ and the radii
$\rho_0,\rho_1\in[R/2,3R/4]$.  Since $M\le2^{p/2}h_0$,
\[
\lambda_0^{(2)}(-h_0)+(1-\lambda_0^{(2)})M
\le h_0\bigl(-\lambda_0^{(2)}+(1-\lambda_0^{(2)})2^{p/2}\bigr)
= h_0\cdot\frac{-(2^{p/2}+\tfrac12)+\tfrac12 2^{p/2}}{2^{p/2}+1}
=-\frac{h_0}2 ,
\]
so those levels lie in $[-h_0,-h_0/2]$, which is inside the range $h\ge-h_0$ on
which (4c) declares (3.12$^+$).  For such a level $l$ and such radii,
(3.12$^+$) at $h=l$, $\rho=\rho_1$ and $\sigma=(\rho_1-\rho_0)/\rho_1$ --- so
that $\rho-\sigma\rho=\rho_0$ and $\sigma\rho=\rho_1-\rho_0$ --- gives
\[
\int_{B_{\rho_0}\cap\{v>l\}}\|V\|^2
\le C_{12}\frac{h_0^2}{(\rho_1-\rho_0)^2}\bigl|B_{\rho_1}\cap\{v>l\}\bigr| ,
\]
and in this branch $h_0^2\le16\operatorname{ess\,sup}_{B_{\rho_1}}\bigl((v-l)^+\bigr)^2$.
Indeed $(v-l)^+\ge v-l$ pointwise, and the essential supremum is translation
equivariant and monotone in the domain, so, using $B_{R/2}\subseteq B_{\rho_1}$
(valid because $\rho_1\ge R/2$),
\[
\operatorname{ess\,sup}_{B_{\rho_1}}(v-l)^+
\ \ge\ \operatorname{ess\,sup}_{B_{R/2}}(v-l)
\ =\ M_k-l\ \ge\ -\tfrac{h_0}4+\tfrac{h_0}2=\tfrac{h_0}4>0 ,
\]
whence $h_0^2\le16\bigl(\operatorname{ess\,sup}_{B_{\rho_1}}(v-l)^+\bigr)^2$ and,
the quantity being nonnegative, $h_0^2\le16\operatorname{ess\,sup}_{B_{\rho_1}}
\bigl((v-l)^+\bigr)^2$.  Together with the essential boundedness of $(v-l)^+$
this is
$\mathrm{DeGiorgiLevelBound}\ n\ (16C_{12})\ v\ V\ x_0\ l\ \rho_0\ \rho_1$
(\noderef{DeGiorgiLevelBound}), which is the instance Step 5's choice of
$\gamma_\sharp$ accounts for.  If $M\le\kappa$ then
$v\le-h_0$ a.e.\ on $B_S\supseteq B_{R/2}$, hence $g^k\le-\lambda$ and
$|g^k|\ge\mu/2$ a.e.\ there and we are done; otherwise $\kappa<M$ and the node's
conclusion is $v\le-h_0/2$ a.e.\ on $B_{\lambda_1'S}=B_{R/2}$, i.e.\ $g^k<0$ and
$|g^k|^{p/2}\ge\tfrac12\lambda^{p/2}$, so
\[
|g^k|\ \ge\ 2^{-2/p}\lambda\ =\ 2^{-1-2/p}\mu\ \ge\ 2^{-1-4/p}\mu
\ \ge\ \min\{\tfrac1{16},2^{-1-4/p}\}\mu
\qquad\text{a.e.\ on }B_{R/2},
\]
using $2/p\le4/p$.  This completes part (a).

\emph{Step 7: part \textup{(b)}.}  Assume the hypothesis of (b): for every $k$,
\[
\bigl|B_R\cap\{g^k>\tfrac\mu2\}\bigr|<\theta_2|B_R|
\quad\text{and}\quad
\bigl|B_R\cap\{g^k<-\tfrac\mu2\}\bigr|<\theta_2|B_R| .
\]
Fix $k$; we prove $g^k\le\nu\mu$ a.e.\ on $B_{R/2}$, the bound
$g^k\ge-\nu\mu$ following by the same argument applied to $-g^k$, whose
hypothesis is the second displayed inequality.

\textbf{(7a) The inequalities (3.21) of the thesis and (3.22) of the thesis.}  With
$\lambda_1S_b=tR$ and $|B_{tR}|=t^n|B_R|$,
\[
\bigl|B_{\lambda_1S_b}\cap\{g^k>\tfrac\mu2\}\bigr|
\le\bigl|B_R\cap\{g^k>\tfrac\mu2\}\bigr|
<\theta_2|B_R|=\frac{\theta_2}{t^n}\bigl|B_{\lambda_1S_b}\bigr|
=c_0\bigl|B_{\lambda_1S_b}\bigr| ,
\tag{3.21}
\]
and likewise with $\{g^k<-\mu/2\}$, which is (3.22) of the thesis.  The constant
$c_0=2\theta_2/(1+\theta_2)$ is $<1$ because $\theta_2<1$, as
\noderef{LevelSetMeasureImprovement} requires.

\textbf{(7b) The measure improvement (3.24) of the thesis.}  Apply
\noderef{LevelSetMeasureImprovement} --- the instance of Step 5 --- with the
centre $x_0$, the radius $S_b$, the level $\kappa=\mu/2$, the function $f=g^k$,
the field $Dg^k$ and $M=\operatorname{ess\,sup}_{B_{S_b}}g^k$.  Its hypotheses:
$0<S_b$; $g^k\in W^{1,2}(B_{S_b})$ with weak gradient $Dg^k$ by (0.3), since
$S_b<R$; $g^k\le\mu$ a.e.; the measure hypothesis is (3.21) of the thesis, whose
ball is $B_{\lambda_1S_b}$; and the family of $(2.1)$ instances is required for
the levels $l\in[\kappa,M]\subseteq[\mu/2,\mu]$ at the single pair of radii
$(\lambda_1S_b,S_b)=(tR,\tfrac{1+t}2R)$, with $tR<\tfrac{1+t}2R<R$ because
$t<1$.  Those are exactly the levels covered by Step 2(ii) when
$\varepsilon<\lambda^2$, and by Step 2(i) when $\varepsilon\ge\lambda^2$, so
Step 3 supplies them with $\gamma=\gamma_\sharp$ in either case.  (This is the
only point at which part (b) needs (3.23) of the thesis rather than (3.9) of the thesis,
and it is why Step 2(ii) was proved in the level-set form.)  If $M\le\kappa$
then $g^k\le\mu/2\le\nu\mu$ a.e.\ on $B_{S_b}\supseteq B_{R/2}$ --- recall
$\nu=1-\lambda_2/4\ge3/4$ --- and we are done; so assume $\kappa<M$.  The
conclusion is
\[
\bigl|B_{\lambda_1S_b}\cap\{\lambda_2\tfrac\mu2+(1-\lambda_2)M<g^k\}\bigr|
\le\theta_1\bigl|B_{\lambda_1S_b}\bigr| .
\tag{3.24}
\]
Since $M\le\mu$, the level satisfies
$\lambda_2\tfrac\mu2+(1-\lambda_2)M\le\mu(1-\tfrac{\lambda_2}2)=:\kappa_2$, so
$\{g^k>\kappa_2\}$ is contained in the set measured in (3.24) of the thesis and
\[
\bigl|B_{\lambda_1S_b}\cap\{\kappa_2<g^k\}\bigr|
\le\theta_1\bigl|B_{\lambda_1S_b}\bigr| .
\tag{3.24$'$}
\]

\textbf{(7c) The maximum principle on the inner pair of balls.}  Apply
\noderef{LevelTruncationMaximumPrinciple} --- the instance of Step 5 at
$\lambda_0=\tfrac12$ and $\lambda_1=\lambda_1''$ --- with the centre $x_0$, the
radius $\lambda_1S_b=tR$, the level $\kappa_2$, the function $f=g^k$, the field
$Dg^k$ and $M'=\operatorname{ess\,sup}_{B_{tR}}g^k$.  Its hypotheses:
$g^k\in W^{1,2}(B_{tR})$ with weak gradient $Dg^k$ by (0.3), since $tR<R$;
$g^k\le\mu$ a.e.; the measure hypothesis is the display (3.24$'$), since
$\theta_1\le\theta^{(3)}$; and the $(2.1)$ instances are required for the levels
$l\in[\kappa_2,\tfrac12\kappa_2+\tfrac12M']$ and the radii
$\rho_0,\rho_1\in[\lambda_1''\cdot tR,\,tR]=[R/2,tR]$ with $\rho_0<\rho_1<R$.
Every such level satisfies $l\ge\kappa_2=\mu(1-\lambda_2/2)\ge\mu/2$ and
$l\le M'\le\mu$, so it lies in $[\mu/2,\mu]$ and Step 2 together with Step 3
supplies the instance, with $\gamma=\gamma_\sharp$.  If $M'\le\kappa_2$ then
$g^k\le\kappa_2=\mu(1-\lambda_2/2)\le\mu(1-\lambda_2/4)=\nu\mu$ a.e.\ on
$B_{tR}\supseteq B_{R/2}$ and we are done; otherwise $\kappa_2<M'$ and the
conclusion is
\[
g^k\le\tfrac12\kappa_2+\tfrac12M'
\le\tfrac12\mu\Bigl(1-\frac{\lambda_2}2\Bigr)+\tfrac12\mu
=\mu\Bigl(1-\frac{\lambda_2}4\Bigr)=\nu\mu
\qquad\text{a.e.\ on }B_{\lambda_1''\cdot tR}=B_{R/2},
\]
using $M'\le\mu$.  Applying the same argument to $-g^k$ --- whose weak gradient
is $-Dg^k$ by \noderef{WeakGradientLinear}, whose essential supremum on each of
the balls above is at most $\mu$, and for which (3.22) of the thesis plays the role of
(3.21) of the thesis --- gives $-g^k\le\nu\mu$ a.e.\ on $B_{R/2}$.  Hence
$|g^k|\le\nu\mu$ a.e.\ on $B_{R/2}$ for every $k$, which is (b).

\medskip
\noindent\textbf{Correction reported.}  Following the paper's Note for the
formalization, the constant $1/16$ in part (a) is replaced by
$\min\{1/16,\ 2^{-1-4/p}\}$.  This equals $1/16$ for $p\ge4/3$; the step
$-\lambda/4^{2/p}<-\lambda/8$ on thesis p.~24 --- Branch 1 of (6b) above ---
requires $4^{2/p}<8$, i.e.\ $p>4/3$, so for $1<p\le4/3$ the constant must be
weakened to $2^{-1-4/p}=4^{-2/p}/2$.  Every branch of part (a) above delivers at
least $\min\{1/16,2^{-1-4/p}\}\mu$: (6a) gives $\mu/8$, Branch 1 of (6b) gives
$2^{-1-4/p}\mu$ and Branch 2 gives $2^{-1-2/p}\mu\ge2^{-1-4/p}\mu$, and the two
degenerate sub-cases give $\mu/2$.
\end{proof}
